% Samanliknande tabell over undergreinene (for kapittelet «Greinene»).
% Kvar grein har ein prosess og eit vokabular, men inga prøvbar modell av forma.
\begin{table}[t]
\centering
\small
\begin{tabular}{@{}p{24mm}p{30mm}p{28mm}p{34mm}@{}}
\toprule
\textsc{grein} & \textsc{prosess} & \textsc{vokabular} & \textsc{det som manglar} \\
\midrule
Industri\-design & styling, ergonomi, produktutvikling & form, funksjon, brukar & ein modell av kva ei god form er (unnataket: kroppen, i ergonomien) \\[4pt]
System\-orientert & gigamapping, boundary critique & node, relasjon, kompleksitet & ein storleik å optimere mot; ein stoppregel; eit falsifiseringsvilkår \\[4pt]
Teneste\-design & blueprint, kundereise & touchpoint, frontstage & ein dom om kva ei teneste \emph{bør} gjere; eit nei \\[4pt]
Interaksjon / \textsc{ux} & brukartesting, heuristikk & affordans, usability & eit kriterium for om noko \emph{bør} lagast, ikkje berre om det verkar \\
\bottomrule
\end{tabular}
\caption{\textsc{Det same fråvêret, fire gonger.} Kvar grein av faget har ein prosess og eit vokabular, men ingen har ein delt, prøvbar modell av kva forma er og kva som verkar på henne. Greinene multipliserte fråvêret; dei fylte det ikkje.}
\end{table}
